\section{Introduction}
Multi-object tracking (MOT) links object detections across time while preserving identities. In modern online trackers, this process is dominated by \emph{association}: motion cues predict where a target should be, and appearance cues decide which detection still belongs to the same identity. While recent tracking-by-detection systems are already strong \cite{zhang2021bytetrack,cao2022ocsort,aharon2022botsort,du2022strongsort}, crowded scenes still expose a fundamental weakness: appearance similarity is helpful precisely when it is also most unreliable.

This failure mode is easy to observe in practice. During long occlusion or dense crowd interaction, a track's most recent ReID feature may become stale, noisy, or contaminated by mismatched updates. If the tracker trusts this instantaneous appearance too much, identity switches follow. If it ignores appearance altogether, association collapses when motion is ambiguous. The key question is therefore not whether appearance should be used, but \emph{when appearance should be trusted}.

Most current MOT systems answer this question implicitly through fixed thresholds, heuristic weighting, or architecture-specific refinements. Our goal is to answer it explicitly with a small, interpretable, and transferable module. To this end, we propose \textbf{Laplace-Guided Temporal Reliability Association (LTRA)}, a plug-in association mechanism that augments a strong baseline tracker with temporal reliability estimation. Instead of learning a new detector, backbone, or end-to-end tracking system, LTRA uses a track's history of appearance features to build \emph{multi-scale Laplace temporal signatures}. These signatures summarize how a target looks over time using exponentially decayed prototypes at several temporal scales. LTRA then measures how well the current detection agrees with those signatures, combines this signal with trajectory stability and motion consistency, and uses the resulting reliability to decide how much the tracker should trust appearance for the current pairwise match.

This design is intentionally narrow. We do not claim a new detector, a new ReID backbone, or a new global tracking framework. We keep the baseline tracker intact and change only the association logic. This makes the method easy to interpret, easier to evaluate fairly, and suitable for a clean paper story: \emph{a temporal reliability estimator that tells an existing tracker when to trust appearance}.

Our implementation targets a BoT-SORT style tracker as the main carrier because it already provides a strong and well-understood combination of motion, appearance, and camera-motion compensation \cite{aharon2022botsort}. On top of that baseline, LTRA acts as a light-weight plug-in in the primary appearance-aware association stage. The central question is narrow and testable: \emph{can a small Laplace-inspired trust calibrator make appearance matching more trustworthy in crowded MOT without changing the detector or ReID backbone?}

\textbf{Contributions.}
\begin{enumerate}
  \item We propose \textbf{Laplace temporal signatures}, a multi-scale exponentially decayed summary of track appearance history that captures both short-term identity evidence and longer-term temporal consistency.
  \item We introduce \textbf{Laplace-Guided Temporal Reliability Association (LTRA)}, which estimates pairwise reliability from temporal agreement, track stability, and motion consistency, and uses it to control how strongly appearance cues affect the final association score.
  \item We formulate LTRA as a \emph{plug-in} module for strong standard trackers, enabling fair same-base comparisons that isolate association gains without retraining the detector or redesigning the ReID encoder.
  \item We present a benchmark plan centered on association-sensitive evidence---HOTA, AssA, IDF1, IDSW, and fragmentation---together with transfer and ablation studies that distinguish LTRA from simple history averaging or single-scale smoothing.
\end{enumerate}
